\chapter{Conclusion}\label{ch:conclusion}

This work addressed the challenge of repetitive manual configuration during 3MF file import in slicer software. Through conceptual exploration and partial implementation, it demonstrated that metadata embedded in a 3D model — specifically, appearance information — can be leveraged to streamline the import process and improve parameter consistency.

The research and experimentation presented throughout Chapters 3 and 4 confirm the technical feasibility of this approach, even though further development is required for full integration. The 3MF format, with its open, extensible, and human-readable structure, proved particularly well-suited to this goal. Its capabilities remain underutilised in current additive manufacturing workflows, representing an opportunity for innovation in user-oriented print preparation tools.

While the implemented prototype represents only a preliminary stage, it establishes a foundation for a more user-friendly and efficient workflow between modelling and slicing environments. A wider adoption of the 3MF format’s capabilities could foster a more standardised exchange of print-related data and encourage community-driven enhancement of open-source slicer tools.

The proposed print hints concept demonstrates a practical method for integrating usability improvements directly into existing workflows. However, the ongoing evolution of slicer technologies and the diversity of user requirements leave ample room for alternative or complementary approaches. Future contributions from the community could further refine or expand upon the ideas presented in this work, ultimately contributing to a more seamless and accessible additive manufacturing process.